\section{Making a Large Island}
\label{sec:graph-making-large-island}

\problemstatement{Given an $n \times n$ binary grid, change \textbf{at
most one} cell from $0$ to $1$ to maximize the size of the largest
resulting island ($4$-connected group of $1$s). Return that maximum
possible size.}

\begin{patternbox}
This closing problem of the chapter combines Disjoint Set's component
\textbf{sizing} (already used implicitly by union-by-size throughout
this chapter) with a final scan: first, group every existing island
using Disjoint Set, tracking each group's size; then, for every water
cell, compute what size would result from flipping \textbf{just that
one cell} -- the sum of the sizes of its (up to $4$) \textbf{distinct}
neighboring islands, plus $1$ for the flipped cell itself.
\end{patternbox}

\subsection*{Why distinctness of neighboring islands matters}

A water cell might have two or more $4$-directional neighbors that
belong to the \emph{same} island (e.g., a water cell with land both
above and below, where that land happens to already be connected via
some other path around the grid). Counting that island's size twice
would overcount. So neighboring islands must be collected into a
\textbf{set} of distinct roots before summing their sizes -- exactly
the kind of deduplication Disjoint Set's $\textit{find}$ operation
makes trivial (two neighbors are ``the same island'' precisely when
$\textit{find}$ returns the same root for both).

\subsection*{Algorithm}

\begin{enumerate}
  \item Build a Disjoint Set over all $n^2$ cells; union every land
        cell with its land neighbors.
  \item For every water cell, collect the distinct roots among its up
        to $4$ neighbors that are land; the candidate size if flipped
        is $1$ plus the sum of those distinct roots' group sizes.
  \item Track the maximum candidate size across all water cells (and
        separately, the largest existing island's size, in case the
        grid has no water cell to flip at all -- the answer is then
        simply $n^2$).
\end{enumerate}

\subsection*{C++ implementation}

\begin{lstlisting}
#include <bits/stdc++.h>
using namespace std;

class DisjointSet {
    vector<int> parent, size;
public:
    DisjointSet(int n) : parent(n), size(n, 1) {
        for (int i = 0; i < n; i++) parent[i] = i;
    }
    int find(int x) {
        if (parent[x] == x) return x;
        return parent[x] = find(parent[x]);
    }
    void unite(int x, int y) {
        int rootX = find(x), rootY = find(y);
        if (rootX == rootY) return;
        if (size[rootX] < size[rootY]) swap(rootX, rootY);
        parent[rootY] = rootX;
        size[rootX] += size[rootY];
    }
    int getSize(int x) { return size[find(x)]; }
};

int largestIsland(vector<vector<int>>& grid) {
    int n = grid.size();
    DisjointSet ds(n * n);
    int dr[] = {-1, 1, 0, 0};
    int dc[] = {0, 0, -1, 1};

    for (int r = 0; r < n; r++) {
        for (int c = 0; c < n; c++) {
            if (grid[r][c] == 1) {
                for (int dir = 0; dir < 4; dir++) {
                    int nr = r + dr[dir], nc = c + dc[dir];
                    if (nr >= 0 && nr < n && nc >= 0 && nc < n && grid[nr][nc] == 1) {
                        ds.unite(r * n + c, nr * n + nc);
                    }
                }
            }
        }
    }

    int best = 0;
    bool hasWater = false;

    for (int r = 0; r < n; r++) {
        for (int c = 0; c < n; c++) {
            if (grid[r][c] == 0) {
                hasWater = true;
                unordered_set<int> distinctRoots;
                for (int dir = 0; dir < 4; dir++) {
                    int nr = r + dr[dir], nc = c + dc[dir];
                    if (nr >= 0 && nr < n && nc >= 0 && nc < n && grid[nr][nc] == 1) {
                        distinctRoots.insert(ds.find(nr * n + nc));
                    }
                }
                int candidate = 1;
                for (int root : distinctRoots) candidate += ds.getSize(root);
                best = max(best, candidate);
            } else {
                best = max(best, ds.getSize(r * n + c));
            }
        }
    }
    return hasWater ? best : n * n;
}

int main() {
    vector<vector<int>> grid = {
        {1, 0},
        {0, 1}
    };
    cout << largestIsland(grid) << endl; // 3
    return 0;
}
\end{lstlisting}

\subsection*{Dry run}

\dryrun{Take the $2\times2$ grid $[[1,0],[0,1]]$ (two separate
single-cell islands, diagonally placed).}

After building the Disjoint Set: cell $(0,0)$ and cell $(1,1)$ are
each their own island of size $1$ (they are not $4$-directionally
adjacent to each other). Checking water cell $(0,1)$: neighbors
$(0,0)$ (land, size $1$) and $(1,1)$ (land, size $1$) -- two
\emph{distinct} roots, candidate $=1+1+1=3$. Checking water cell
$(1,0)$: neighbors $(0,0)$ (land, size $1$) and $(1,1)$ (land, size
$1$) -- also distinct, candidate $=1+1+1=3$. Best: $3$ -- flipping
either water cell connects both existing islands into one.

\begin{complexitybox}
\textbf{Time Complexity.} $O(n^2 \cdot \alpha(n^2))$ -- building the
Disjoint Set and the final scan both touch every cell a constant
number of times.
\textbf{Space Complexity.} $O(n^2)$ for the Disjoint Set.
\end{complexitybox}

\begin{takeawaybox}
This chapter closes by combining nearly everything it introduced:
Disjoint Set's grouping (Section~\ref{sec:graph-disjoint-set-impl}),
size tracking (the same union-by-size bookkeeping Kruskal's and Prim's
relied on internally), and the deduplication that distinguishes a
\emph{group} from a list of individual elements (echoing Accounts
Merge's account-grouping and Most Stones Removed's row/column
grouping) -- a fitting capstone for a chapter built entirely around one
data structure's surprising range of applications.
\end{takeawaybox}
